\documentstyle[11pt]{article}

\pagestyle{empty}

\topmargin -.8in

\textheight 9.5in

\textwidth 6.2in

\title{ An Algorithm And Its Data Structure From Sequential $US \ \ MMCM$

To Parallel Machine}

\author{Hu Guoding \& Wang Yongge \\ Nankai Institute of Mathematics

\\ Nankai University, Tianjin, P.R.China}

\date{}

\begin{document}

\maketitle

\thispagestyle{empty}

\begin{abstract}

 In paper $[1]$,   Hu Guoding gave a general definition of  computing machine model

$MMCM$,   and in $[2]$,  Hu proved that one kind of this machine,   called $US \ \ MMCM$,

can be implemented parallely. In this paper,   we will give an algorithm to

carry out this transformation from sequential $US \ \ MMCM$ to parallel $US \ \ MMCM$,   and

also we will describe the software TR.LSP of this transformation on micro-computer

PC/AT.

\end{abstract}



\section{Computing Machine Model $MMCM$}



\subsection{Notations}



Let $B$ be any set,   $N$ be a natural number, $<B>_N=\{ \alpha |\alpha\in

B^*,   |\alpha|\leq N\}$.



If $v=a_1a_2a_3\cdots a_n,   a_i\in A$,   then $\stackrel{\downarrow\mbox{ } }{v}

=\stackrel{\downarrow}{a}_1a_2a_3\cdots a_n, $ and

$\stackrel{\downarrow\mbox{ }}{v}=\epsilon$ for

$n=0$;  $\stackrel{\mbox{ }\downarrow}{v}=

a_1a_2a_3\cdots\stackrel{\downarrow}{a}_n, $ and

$\stackrel{\mbox{ }\downarrow}{v}=\epsilon$ for

$n=0$;  $\stackrel{\downarrow\downarrow}{v}=

\stackrel{\downarrow}{a}_1a_2a_3\cdots \stackrel{\downarrow}{a}_n, $

and $\stackrel{\downarrow\downarrow}{v}=\epsilon$ for

$n=0$.



\subsection{Computing Machine Model $MMCM$}



For the detailed information of $MMCM$,   the reader is refered to $[1]$,

here we will give a concise description of $MMCM$.



{\bf Definition:} an $MMCM$ is a 6-tuple $(A,   W,   Q,   Q_h,   F,   P)$ where



(1). $A$ is a finite alphabet set;



(2). $W\subset A^*$ is the initial words set;



(3). $Q$ is the finite states set and $Q_h\subset  Q$ is the halting states set;



(4). $F:E_1\times (Q-Q_h)\longrightarrow E_2\times Q$;

where $E_1\subset <A'\cup V>_N,   E_2\subset  <A\cup A'\cup V'>_N$ and

$V=\{ v_1,  v_2,  \cdots ,  v_n\},

V'=\{\stackrel{\downarrow\mbox{ }}{v},

\stackrel{\mbox{ }\downarrow}{v},

\stackrel{\downarrow\downarrow}{v}|v\in V\} ,

A'=\{\stackrel{\downarrow}{a}|a\in A\}$,   the

elements of $E_1$ is in the form of

$\bar{\alpha}=v_1\stackrel{\downarrow}{a}_1v_2\stackrel{\downarrow}{a}_2\cdots v_m\stackrel{\downarrow}{a}_mv_{m+1};

2m+1\leq N $.

If $F(\bar{\alpha},  q)=(\bar{\alpha}_1, q_1)$, then $(\bar{\alpha}_1, q_1)$

is only determined by

$a_1,   a_2,   \cdots,   a_m,   q$; if $v$ or

$\stackrel{\downarrow\mbox{ }}{v},

\stackrel{\mbox{ }\downarrow}{v},   \stackrel{\downarrow\downarrow}{v}$

ocur in $\bar{\alpha}_1,$ then

$v$ must occur in $\bar{\alpha}$.



(5). $P$ is the computation processes set,   e.g. the following is a computation

process:



\centerline{$(\alpha_{\lambda_0},  q_0)\rightarrow

(\alpha_{\lambda_1},  q_1)\rightarrow \cdots\rightarrow (\alpha_{\lambda_t},  q_t)$}.



Now,   we will show how to implement one step of this computation:



\centerline{$(\alpha_{\lambda_i},  q_i)\rightarrow(\alpha_{\lambda_{i+1}},  q_{i+1});  \ \

\alpha_{\lambda_i} \in (A\cup A')^*$}



Firstly,   let $\alpha_{\lambda_i}=a_1\cdots\stackrel{\downarrow}{a}_{i_1}\cdots\stackrel{\downarrow}{a}_{i_m}\cdots a_n$

and $v_1=a_1a_2\cdots a_{i_1-1};  v_2=a_{i_1+1}\cdots a_{i_2-1};   \cdots ;

v_m=a_{i_{m-1}}\cdots a_{i_m-1};  v_{m+1}=a_{i_m+1}\cdots a_n$. Thus,   we have

$\bar{\alpha}_{\lambda_i}=v_1\stackrel{\downarrow}{a}_{i_1}v_2\cdots

v_m\stackrel{\downarrow}{a}_{i_m}v_{m+1}$,   using $F$,   we can

have $F(\bar{\alpha}_{\lambda_i}, q_i)=(\bar{\alpha}_{\lambda_{i}'},   q); $

replace

the $v_i$ in $\bar{\alpha}_{\lambda_i}$ with the above definitions of $v_i$,

we get $\alpha_{\lambda_{i+1}}$;  and $q'=q_{i+1}$.



Lastly,   we need that $q_0$ is the initial state and $q_t\in Q_h$. If

we remove the arrows in $\alpha_{\lambda_0}$,   it becomes an element of $W$.



\section{Uniformly Separable Machine $US \ \ MMCM$}



{\bf Definition:} An $MMCM$  \ $M$ is a Uniformly

Separable $MMCM$,   if its instruction set

$F:E_1\rightarrow E_2$ (or we can write it as

$(v_1\stackrel{\downarrow}{a}_1v_2\cdots

v_m\stackrel{\downarrow}{a}_mv_{m+1}, q)

\stackrel{F}{\rightarrow} (\bar{\alpha},   q)$)

satisfies the following conditions:



(1). the length of $\bar{\alpha}$  only depends on $m,   q$;



(2). $q'$ only depends on $m,   q$;



(3). the letters of $(V\cup V')$ which occur in $\bar{\alpha}$

and their occuring position

in $\bar{\alpha}$ only depend on $m, q$;



(4). the occuring-positions of letters of $A$ in

$\bar{\alpha}$ only depend on $m,   q$;



(5). which letters of $A\cup A'$ will occur in $\bar{\alpha}$

only depends on $a_1, a_2, \cdots, a_m, m$ and  $q$.



\section{The Transformation Algorithm From Sequential

$US \ \ MMCM$ to Parallel $US \ \ MMCM$}



\subsection{the data structure of storing a $US \ \ MMCM$}





Although an $MMCM$ M is a 6-tuple $(A,  W,  Q,  Q_h,  F,  P)$,

we need only to store the $W,   Q_h$ and $F$,   because

$A$ can be a default general finite alphabet

set,   $Q$ belongs to $F$,   $P$ can be computed from

$W, F$ and $Q_h$. Without loss of generality, we can assume

that there is only one element in $Q_h$ and it is stored in

the form of an alphabet string in computer. The computation process

$P$ of a $US \ \ MMCM$ is always the same when the length

of input words and the operating letters coodinates are

the same,   so for an $N$-length input words with the operating

letters coodinates $(\lambda_1\ \  \lambda_2\ \ \cdots\ \ \lambda_s)$,

we can store it in a variable (((a 1)\  (a 2)\  (a 3) \  $\cdots$ \  (a N) \

($\lambda_1$ \  $\lambda_2$\   $\cdots$\    $\lambda_s$)). For the

convenience of LISP language, we use (a i) to denote the

variable $a_i$. Any one production rule of $F$

can be stored in the form of structure(which equivelents

to a record in Pascal language). That is, a Production structure

is made up of the fields $premise,  poperating, pstate, result,

roperating,  rnewborn$ and $rstate$, here $premise$ is the

premise, stored in the form of ((a 1)\  (a 2) (a 3)\ $\cdots$\ (a n));

$poperating$ is the operating letters coodinates, stored in

the form of ($i_1$  $i_2$  $\cdots$ $i_s$);

$result$ is the result which has the same data structure

as $premise$; $roperating$ is operating letters coodinates

for next operation, stored in the form of

$((i_1\ s_1)\ \ (i_2\ s_2)\ \  \cdots\ \  (i_t \ s_t))$,

where $i_j$ is the operating letters coodinates and

$s_j=0,\ 1,\ 2$ or $3$, here $(i_j\  0), \  (i_j\  1), \

(i_j \ 2)$ and $(i_j \ 3)$ denote

$v, \stackrel{\downarrow\mbox{ }}{v}, \stackrel{\mbox{ }\downarrow}{v}

and \stackrel{\downarrow\downarrow}{v}$ respectively;

$rnewborn$ is newborn words, stored in the form of

$(\mu_1\ \mu_2\ \cdots\  \mu_t)$; $pstate$ and $ratate$

are states, stored in the form of alphabet string. After

the above structural  process,   the whole $F$ can then be

stored in an array $proarry$ which has structure as its components.



\subsection{the algorithm to obtain $US \ \ MMCM$'s computation process}



Generally, production rule is in the form of

$v_1c_1v_2c_2\cdots v_nc_nv_{n+1}\rightarrow v_{i_1}b_1v_{i_2}b_2\cdots v_{i_m}b_mv_{i_{m+1}}$,

where $v_i$ are string variables and $c_i, b_i$ are letters.

because the initial word is $a_1a_2\cdots a_n$,

to obtain the computation, we must replace

subwords between operating letters with the variables, i.e.

if $\alpha_0=a_1\cdots\stackrel{\downarrow}{a}_{i_1}\cdots\stackrel{\downarrow}{a}_{i_m}\cdots a_n$,

then let $v_1=a_1a_2\cdots a_{i_1-1}; v_2=a_{i_1+1}\cdots a_{i_2-1}; \cdots ;

v_m=a_{i_{m-1}}\cdots a_{i_m-1}; v_{m+1}=a_{i_m+1}\cdots a_n$,

so we have

$\alpha_0'=v_1\stackrel{\downarrow}{a}_{i_1}v_2\cdots v_m\stackrel{\downarrow}{a}_{i_m}v_{m+1}$.

After doing this, we use the appropriate

production rule in $F$ to carry out the computation,

and then, replace the $v_i$'s with the old value. This

computation process is the task of function $adapt$ and

$huanyan$ in the software {\bf TR.LSP}. The complexities

of $adapt$ and $huanyan$ are $O(n)$ and $O(n^2)$ respectively.

From above analysis, each one step of the computation process

$P$ can be obtained from the following computations:



(1). $adapt$;



(2). search for the approprate production rule(which is

completed by the function $search$);



(3). $huanyan$.



\subsection{find out the maximal parallelism in $P$}



Step 1: find the $PSCPP$'s in $P$(Pure Serial Computation

Partial Process), i.e. the $k_1, k_2, \cdots, k_t$ in paper

$[2]$. To do this, for each step of $P$, we test whether

there are any operating letters of this step which belong to the

newborn words of last step. The result of this step is

stored in the array $klevel$, i.e. $klevel[i]$ is made up of

the steps of $P$ which are in $k_i$. the complexity of this

step is $O(n)$.



\noindent

Step 2: find out the maximum parallelism for each $k_i$ in $PSCPP$,

i.e. we classify $PSCPP$ into $I_1, I_2, \cdots, I_s$

so that for each $I_j$, the elements $k_i$ in $I_j$ can be

implemented parallely. To get $I_j$, firstly, we must have the

operating letters for each $k_i$(Remark: $k_i$ is a composite

map of simple $US \ \ MMCM$ map, its operating letters is

the union of the submaps's operating letters) and these

operating letters are produced by which $k_s$, secondly,

using the above result to label each $k_i$, e.g. if

the operating letters of $k_i$ are produced by

$k_{j_1}, k_{j_2}, \cdots, k_{j_t}$,

then give $k_i$ the label $\max(j_1, j_2, \cdots, j_t)$.

Lastly, using the labels to get the classification. The

complexity of step 2 is $O(n^3)$.



\section{The Auto-transformation Software From

Sequential Algorithm to Parallel Algorithm}



We have designed a software {\bf TR.LSP} which can be

used on PC/AT to carry out this transformation. The

function of this software is: iuput a sequential

$US \ \ MMCM$, then the software can output the parallel

machine equivalent to the inputed sequential machine. I.e.

output the computation process $P$:



\centerline{$\alpha_0 \stackrel{1*}{\longrightarrow}\alpha_1\stackrel{2*}{\longrightarrow}\alpha_2\stackrel{3*}{\longrightarrow}\cdots\stackrel{n*}{\longrightarrow}\alpha_n$}



\noindent

and classify this process into: $I_1,  I_2,  \cdots,   I_l$ which satisfies:



(1). $\cup_{i=1}^lI_i=\{ 1*,   2*,  \cdots,   n*\} $;



(2). $I_i\cap I_j=\phi,   i\not= j$;



(3). If $I_i=\{ n_1*,   n_2*,   \cdots,   n_k*\}$,   then,   in the process $P$,

the $n_1*,   n_2*,   \cdots,   n_k*$-st step can be implemented parallely.



In other words,   the output parallel machine operates as follows:



Step 1. carry out each steps in $I_1$ parallely;



Step 2. carry out each steps in $I_2$ parallely;



\ \ \ \ \ \  \  \ $\cdots$  \ \ \ \ \ $\cdots$



Step l. carry out each steps in $I_l$ parallely;

and output the whole result.



The strucral scheme of this software can be found

at the end of this article.



\section{An application of This Software}



Example: one application of software TR.LSP: Mulplication of numbers:



\vskip 3pt

\noindent

(I). the production rules for $US \ \ MMCM$  M of mulplication



( ((v 1) (a 1) (a 2) (a 3)) (2 3 4) ``q0" ((v 1) (b 1) (b 2) (a 2) (a 3) (b 3)

(b 4)) ((1 2) (4 0) (5 0)) (2 3 6 7) ``q1")



( ((v 1) (a 1) (v 2) (a 2) (a 3) (v 3)) (2 4 5) ``q1" ((v 1) (b 1) (b 2) (v 2) (a 2)

(a 3) (b 3) (b 4) (v 3)) ((1 2) (5 0) (6 0)) (2  3 7 8) ``q1")



( ((a 1) (v 1) (a 2) (a 3) (v 2)) (1 3 4) ``q1" ((b 1) (b 2) (v 1) (a 3) (b 3)

(b 4) (v 2)) ((3 2) (4 0)) (1 2 5 6) ``q1")



( ((v 1) (a 1) (a 2) (v 2)) (2 3) ``q1" ((v 1) (a 1) (a 2) (v 2)) ((1 2) (3 0))

nil ``q1")



( ((v 1) (a 1) (v 2) (a 2) (v 3)) (2 4) ``q1" ((v 1) (a 1) (v 2) (a 2) (v 3))

((1 2) (2 0) (4 0)) nil ``q2")



( ((v 1) (a 1) (v 2) (a 2) (v 3)) (2 4) ``q2" ((v 1) (a 1) (v 2) (a 2) (v 3))

((1 2) (2 0) (4 0)) nil ``q2")



( ((v 1) (a 1) (a 2) (v 2) (a 3) (v 3)) (2 3 5) ``q2" ((v 1) (b 1) (b 2) (v 2)

(a 3) (v 3)) ((1 2) (5 0)) (2 3) ``q2" )



( ((a 1) (v 1) (a 2) (v 2)) (1 3) ``q2" ((a 1) (v 1) (a 2) (v 2)) ((2 2) (3 0))

 nil ``q3" )



( ((a 1) (a 2) (v 1) (a 3) (v 2)) (1 2 4) ``q2" ((b 1) (b 2) (v 1) (a 3) (v 2))

((3 2) (4 0)) (1 2) ``q3" )



( ((v 1) (a 1) (a 2) (v 2)) (2 3) ``q3" ((v 1) (a 1) (a 2) (v 2)) ((1 2) (3 0))

 nil ``q3" )



( ((v 1) (a 1) (v 2) (a 2) (v 3)) (2 4) ``q3" ((v 1) (a 1) (v 2) (a 2) (v 3))

((1 2) (4 0)) nil  ``q4" )



( ((v 1) (a 1) (v 2) (a 2) (v 3)) (2 4) ``q4" ((v 1) (a 1) (v 2) (a 2) (v 3))

((1 2) (4 0)) nil ``q5" )



( ((v 1) (a 1) (v 2) (a 2) (v 3)) (2 4) ``q5" ((v 1) (a 1) (v 2) (a 2) (v 3))

((1 2) (2 0) (4 0)) nil ``q5" )



( ((v 1) (a 1) (a 2) (v 2) (a 3) (v 3)) (2 3 5) ``q5" ((v 1) (b 1) (v 2)

(a 3) (v 3)) ((1 2) (4 0)) (2) ``q5" )



( ((a 1) (v 1) (a 2) (v 2)) (1 3) ``q5" ((b 1) (v 1) (a 2) (v 2)) ((3 0) (4 2))

(1) ``q7" )



( ((a 1) (a 2) (v 1) (a 3) (v 2)) (1 2 4) ``q5" ((b 11) (v 1) (a 3)

(v 2)) ((3 0) (4 2)) (1) ``q7" )



( ((v 1) (a 1) (v 2) (a 2)) (2 4) ``q7" ((v 1) (a 1) (v 2) (a 2))

 ((2 0) (3 2)) nil ``q7" )



( ((v 1) (a 1) (v 2) (a 2) (v 3)) (2 4) ``q7" ((v 1) (a 1) (v 2) (a 2) (v 3))

((2 0) (3 2) (4 0)) nil ``q8" )



( ((v 1) (a 1) (v 2) (a 2) (v 3)) (2 4) ``q8" ((v 1) (a 1) (v 2) (a 2) (v 3))

((2 0) (3 2) (4 0)) nil ``q8")



( ((v 1) (a 1) (v 2) (a 2) (a 3) (v 3)) (2 4 5) ``q8" ((v 1) (a 1) (v 2) (b 7)

 (b 8) (v 3)) ((2 0) (3 2)) (4 5) ``q8" )



( ((v 1) (a 1) (a 2) (v 2)) (2 3) ``q8" ((v 1) (a 1) (a 2) (v 2)) ((2 0) (4 2))

 nil ``q9" )
�( ((v 1) (a 1) (a 2) (a 3) (v 2)) (2 3 4) ``q8" ((v 1) (a 1) (b 1) (b 2) (v 2))

 ((2 0) (5 2)) (3 4) ``q9" )



( ((v 1) (a 1) (v 2) (a 2)) (2 4) ``q9" ((v 1) (a 1) (v 2) (a 2)) ((2 0) (3 2))

 nil ``q9" )



( ((v 1) (a 1) (v 2) (a 2) (v 3)) (2 4) ``q9" ((v 1) (a 1) (v 2) (a 2) (v 3))

((2 0) (3 2)) nil  ``q10" )



( ((v 1) (a 1) (v 2) (a 2) (v 3)) (2 4) ``q10" ((v 1) (a 1) (v 2) (a 2) (v 3))

 ((2 0) (3 2)) nil ``q11" )



( ((v 1) (a 1) (v 2) (a 2) (v 3)) (2 4) ``q11" ((v 1) (a 1) (v 2) (a 2) (v 3))

 ((2 0) (3 2) (4 0)) nil ``q11" )



( ((v 1) (a 1) (v 2) (a 3) (a 4) (v 3)) (2 4 5) ``q11" ((v 1) (a 1) (v 2) (b 9)

(v 3)) ((2 0) (3 2)) (5) ``q11" )



( ((v 1) (a 1) (a 2) (v 2)) (2 3) ``q11" ((v 1) (a 1) (b 1) (v 2))

((1 2) (2 0) (4 2)) (3) ``q12" )



( ((v 1) (a 1) (a 2) (a 3) (v 2)) (2 3 4) ``q11" ((v 1) (a 1) (b 11)

(v 2)) ((1 2) (2 0) (4 2)) (3) ``q12" )



( ((v 1) (a 1) (a 2) (v 2) (a 3)) (2 3 5) ``q12" ((v 1) (a 1) (a 2) (v 2) (a 3))

((1 2) (3 0) (5 0)) nil ``q13")



( ((v 1) (a 1) (v 2) (a 2) (v 3) (a 3)) (2 4 6) ``q13" ((v 1) (b 1) (b 2) (v 2)

(a 2) (v 3)) ((1 2) (5 0) (6 2)) (2 3) ``q13")



( ((a 1) (v 1) (a 2) (v 2) (a 3)) (1 3 5) ``q13" ((b 1) (b 2) (v 1) (a 2) (v 2))

((4 0) (5 2)) (1 2) ``q13")



( ((v 1) (a 1) (v 2) (a 2)) (2 4) ``q13" ((b 1) (b 2) (v 1) (a 1) (v 2))

((4 0) (5 2)) (1 2) ``q13")



( ((v 1) (a 1) (a 2)) (2 3) ``q13" ((b 1) (b 2) (v 1) (a 1)) ((3 2) (4 0))

(1 2) ``q14")



( ((v 1) (a 1) (a 2)) (2 3) ``q14" ((v 1) (a 1)) ((1 2)) nil ``q14")

( ((v 1) (a 1) (v 2)) (2) ``q14" ((v 1) (a 1) (v 2)) ((1 2)) nil ``q15")



( ((v 1) (a 1) (v 2)) (2) ``q15" ((v 1) (a 1) (v 2)) ((1 2) (2 0)) nil ``q15")

( ((v 1) (a 1) (a 2) (v 2)) (2 3) ``q15" ((v 1) (b 1) (v 2)) ((1 2)) (2) ``q15")



( ((a 1) (a 2) (v 1)) (1 2) ``q15" ((b 1) (v 1)) nil (1) ``qe")

( ((a 1) (v 1)) (1) ``q15" ((b 1) (v 1)) nil (1) ``qe")



nil



\noindent

(II). the computation process for above $MMCM$ when

input length ``2''$\times$``10'' is:



\centerline{$\alpha_0\stackrel{1*}{\longrightarrow}\alpha_1\stackrel{2*}{\longrightarrow}\cdots\stackrel{128*}{\longrightarrow}\alpha_{128}$}

the exact description of $\alpha_i$ are omited here.



\noindent

(III). the parallel machine for multiplication of ``10''$\times$``2'' is:





$I_1=(1 \  2\  3 \  4 \  5 \  6 \ 7 \ 8 \ 9 \ 10)$

 

$I_2=$(11  12  13  14  15  16  17  18  19  20  21  22  23  24  25  26  27  28  29  30  31  50

51  52  53  54  55  56  57  58  59  60  61  62  63  64  65  66  67  68  69  70  71  90  91  104)



$I_3=$(32  33  34  35  36  37  38   39  40  41  42  43  44  45  46  47  48  49  72  73  74

  75  76  77  78  79  80  81  82  83  84  85  86  87  88  89  92  93  94  102  103)



$I_4=$(95 96 97 98 99 100 101 105 106 107 108 109 110 125 126 127)



$I_5$=(111  112  113  114  115  116  117  118  119  120  121  122  123  124  128)



The above result shows that for the multiplication of form ``10''$\times$``2'',

the sequential machine need 128 steps,   but the parallel machine needs only 5

steps: $I_1,   I_2,   I_3,   I_4,   I_5$.



\vskip 10pt



\begin{center}



\begin{tabular}{|l|}\hline

{\em Input the $US \ \ MMCM$ machine}\\

{\em and the interface process}\\ \hline

\end{tabular}



$\downarrow$



\begin{tabular}{|l|}\hline

{\em Adapt the initial words}\\

{\em $a_1\cdots a_N$ to the production}\\

{\em form $a_1v_1\cdots a_mv_m$,}\\

{\em i.e. the function: $adapt$}\\ \hline

\end{tabular}



$\downarrow$



\begin{tabular}{|l|}\hline

{\em search production rule: i.e.$search$}\\ \hline

\end{tabular}



$\downarrow$



\begin{tabular}{|l|}\hline

{\em variables replacement: $huanyan$}\\ \hline

\end{tabular}



$\downarrow$



\begin{tabular}{|l|}\hline

{\em Obtain the computation}\\

{\em process, i.e. $getprocess$}\\ \hline

\end{tabular}



$\downarrow$



\begin{tabular}{|l|}\hline

{\em Label the $PSCPP$}\\

{\em i.e. $getlevel$}\\ \hline

\end{tabular}



$\downarrow$



\begin{tabular}{|l|}\hline

{\em Classification: $getIs$}\\ \hline

\end{tabular}



$\downarrow$



\begin{tabular}{|l|}\hline

{\em The mainbody of software}\\ \hline

\end{tabular}



$\downarrow$



\begin{tabular}{|l|}\hline

{\em Show the results}\\ \hline

\end{tabular}



\vskip 10pt



{\it Fig.1 the structral scheme of TR.LSP}

\end{center}



\vskip10pt



\centerline{\bf References}



$[1]$. Hu Guo-ding. On Mathematical Model of Computing Machine,

J. of Comp. 1988.



$[2]$. Hu Guo-ding.  Parallel Machine Simulating Serial Computing

Machine. {\it in this volume.}



$[3]$. Gold Hill Computer Company,   GOLDEN COMMON LISP REFERENCE.



$[4]$. L.Silklossy,   LET'S TALK LISP,   1976.

\end{document}



